%% ch06 --- Errors: exceptions as control flow, EAFP, ExceptionGroup and the Result pattern
%%
%% Stub, generated by tools/gen_stubs.py from tools/chapters.json.
%% The brief below is the contract for this chapter: what it must argue, what
%% it must buy the reader, and what it must NOT repeat from the companion
%% volumes. Writing the chapter means deleting the \chapterstub{} block and
%% replacing it with the chapter -- after which this file is never regenerated.
%% If the brief turns out to be wrong once the code has been run, change it in
%% the manifest and record the contradiction in CLAUDE.md under *Resolved
%% questions*.

\chapter{Errors: exceptions as control flow, EAFP, ExceptionGroup and the Result pattern}
\label{ch:errors}

\chapterstub{%
Argument: Python raises where C\# checks, and the cost model is different
\dash{} a \code{try} is nearly free, raising is not free, and both are
normal. EAFP against LBYL; \code{KeyError}, \code{StopIteration} and
\code{AttributeError} as protocol rather than failure; no checked
exceptions, and what stands in for them (a docstring, because a type hint
does NOT carry exceptions); chaining with \code{raise from} against
\code{InnerException}; \code{try}, \code{except}, \code{else},
\code{finally}; custom hierarchies; \code{ExceptionGroup} and \code{except*}
against \code{AggregateException}; \code{contextlib.suppress}; Result and
Either patterns and when they are un-Pythonic; errors at the boundary,
\code{ValidationError} from pydantic and \code{HTTPStatusError} from httpx.
Traps to elicit: a bare \code{except}; catching \code{Exception} to log and
continue; \code{raise e} losing the traceback; a \code{return} inside
\code{finally}. Listings: the same failure handled EAFP and LBYL.
Measurement: experiment E3, the cost of a \code{try} against a check on the
happy and the unhappy path, free. Four exercises. Overlap rule: the
LangChain book's chapter 4 owns cancellation and timeouts, and they are not
here.

\medskip
\textbf{Word budget:} 3000.\\
\textbf{Ships in:} v0.1.\\
\textbf{Measurement:} E3.
}
